\begin{center}
{\large \textbf{АНОТАЦІЯ}}
\end{center}

\textit{Шепілев Д. С.} Методика навчання розробки імерсивних освітніх ресурсів майбутніх учителів інформатики.~-- Кваліфікаційна наукова праця на правах рукопису.

Дисертація на здобуття наукового ступеня доктора філософії за спеціальністю 015~-- Професійна освіта.~-- Криворізький державний педагогічний університет, Кривий Ріг, 2026.\\

Дисертацію присвячено теоретичному обґрунтуванню, розробці та експериментальній перевірці методики навчання розробки імерсивних освітніх ресурсів (ІОР) майбутніх учителів інформатики. Актуальність дослідження зумовлена зростаючою потребою освітньої практики в імерсивних ресурсах (VR, AR, MR, XR), недостатнім рівнем готовності майбутніх учителів інформатики до їх розробки та відсутністю розуміння педагогічних механізмів, що забезпечують ефективність навчання імерсивним технологіям.

Мета дослідження~-- теоретично обґрунтувати та експериментально перевірити методику навчання розробки імерсивних освітніх ресурсів майбутніх учителів інформатики.

Методи дослідження: аналіз і систематизація наукових джерел, порівняльний аналіз освітньо-професійних програм, моделювання, педагогічний експеримент (констатувальний, формувальний, контрольний етапи), анкетування, тестування, спрямований якісний контент-аналіз транскриптів відеолекцій, статистичні методи (критерій $\chi^2$ Пірсона), тріангуляція даних у межах конвергентного паралельного дизайну змішаних методів.

На основі аналізу наукових джерел уточнено сутнісні характеристики імерсивних освітніх ресурсів, виокремлено п'ять сутнісних характеристик ІОР та обґрунтовано 6-компонентну модель їх проєктування. Систематизовано вітчизняний досвід на основі аналізу ОПП 38~ЗВО України за моделлю рівнів зрілості інтеграції ІОР. Когнітивно-афективну модель імерсивного навчання (CAMIL) та модель технолого-педагогічних знань (TPACK) визначено наскрізною теоретичною рамкою дослідження.

Визначено комплекс з 11~дидактичних підходів у 4~кластерах та 8~принципів навчання розробки ІОР. Розроблено структурно-функціональну модель підготовки з п'яти блоків та шести компонентів готовності. Обґрунтовано чотири педагогічні умови та триетапну методику навчання (теоретико-аналітичний, практико-розробницький, інтегративно-рефлексивний етапи).

Проведено педагогічний експеримент ($N = 64$, ЕГ~$= 34$, КГ~$= 30$, 2020--2025) на базі КДПУ. Статистично значущу динаміку підтверджено за всіма критеріями ($\chi^2_{\text{емп}} > 7{,}815$ при $\alpha = 0{,}05$): частка студентів ЕГ з достатнім та високим рівнями зросла від 23,5\% до 70,6\%. Спрямований якісний контент-аналіз 14~транскриптів відеолекцій (1005~сегментів, $\approx$23~години) виявив темпоральну динаміку факторів CAMIL, підтвердив реалізацію всіх педагогічних умов та розкрив домінування стратегії моделювання ($\approx$43\%) з переходом від інструктор-центрованих до студент-центрованих стратегій. Тріангуляція кількісних та якісних даних підтвердила конвергенцію результатів за всіма критеріями.

Наукова новизна: \textit{уперше} визначено ІОР як педагогічний об'єкт у системі підготовки вчителів інформатики; обґрунтовано структурно-функціональну модель підготовки з п'яти блоків; розроблено трикрокову методику навчання WebAR з інтеграцією ML; виявлено темпоральну динаміку факторів CAMIL у процесі навчання розробки ІОР. \textit{Удосконалено} класифікацію ІОР; методику навчання розробки WebAR-застосунків; підходи до валідації методик навчання імерсивних технологій шляхом тріангуляції кількісних та якісних даних. \textit{Набуло подальшого розвитку} наукові положення щодо використання VR/AR-технологій в освіті; підходи до інтеграції ML у WebAR-застосунки; підходи до якісного аналізу педагогічного дискурсу в навчанні імерсивних технологій.

Практичне значення: навчальний посібник з WebAR (10 розділів), курс <<Синтетична реальність в освіті>> у LMS Moodle (18 модулів), лабораторний практикум (10 варіантів), віртуальна хімічна лабораторія з доповненою реальністю.

\textbf{Ключові слова}: імерсивні освітні ресурси, WebAR, доповнена реальність, віртуальна реальність, учитель інформатики, професійна підготовка, педагогічний експеримент, змішані методи, якісний контент-аналіз, CAMIL, TPACK, MindAR, Three.js, A-Frame, машинне навчання, TensorFlow.js.

\clearpage

\begin{center}
{\large \textbf{SUMMARY}}
\end{center}

\textit{Shepilev D. S.} Methodology for Teaching the Development of Immersive Educational Resources for Future Computer Science Teachers.~-- Qualifying scientific work as a manuscript.

Dissertation for the degree of Doctor of Philosophy (PhD) in specialty 015~-- Professional Education.~-- Kryvyi Rih State Pedagogical University, Kryvyi Rih, 2026.\\

The dissertation is devoted to the theoretical substantiation, development, and experimental verification of the methodology for teaching the development of immersive educational resources (IER) for future computer science teachers. The relevance of the research is determined by the growing need for immersive resources (VR, AR, MR, XR) in educational practice, the insufficient level of readiness of future computer science teachers for their development, and the lack of understanding of the pedagogical mechanisms that ensure the effectiveness of immersive technology education.

The aim of the study is to theoretically substantiate and experimentally verify the methodology for teaching the development of immersive educational resources for future computer science teachers.

Research methods: analysis and systematization of scientific sources, comparative analysis of educational programs, modeling, pedagogical experiment (ascertaining, formative, and control stages), questionnaires, testing, directed qualitative content analysis of video lecture transcripts, statistical methods (Pearson's $\chi^2$ criterion), data triangulation within a convergent parallel mixed-methods design.

Based on the analysis of scientific sources, the essential characteristics of immersive educational resources have been clarified, five essential IER characteristics have been identified, and a 6-component IER design model has been substantiated. The domestic experience has been systematized based on the analysis of educational programs of 38~Ukrainian HEIs using an IER integration maturity model. The Cognitive-Affective Model of Immersive Learning (CAMIL) and Technological Pedagogical Content Knowledge (TPACK) model have been established as the overarching theoretical framework.

A set of 11~didactic approaches in 4~clusters and 8~principles for teaching IER development has been identified. A structural-functional model of training comprising five blocks and six readiness components has been developed. Four pedagogical conditions and a three-stage teaching methodology (theoretical-analytical, practical-developmental, and integrative-reflective stages) have been substantiated.

A pedagogical experiment was conducted ($N = 64$, EG~$= 34$, CG~$= 30$, 2020--2025) at KRSPU. Statistically significant dynamics were confirmed across all criteria ($\chi^2_{\text{emp}} > 7.815$ at $\alpha = 0.05$): the proportion of EG students at sufficient and high levels increased from 23.5\% to 70.6\%. Directed qualitative content analysis of 14~video lecture transcripts (1,005~segments, $\approx$23~hours) revealed the temporal dynamics of CAMIL factors, confirmed the implementation of all pedagogical conditions, and identified the dominance of the modeling strategy ($\approx$43\%) with a shift from instructor-centered to student-centered strategies. Triangulation of quantitative and qualitative data confirmed convergence across all criteria.

Scientific novelty: \textit{for the first time}, IER was defined as a pedagogical object in the system of computer science teacher training; a structural-functional model of training comprising five blocks was substantiated; a three-step methodology for teaching WebAR with ML integration was developed and tested; temporal dynamics of CAMIL factors in the IER development training process were identified. \textit{Improved}: classification of IER; methodology for teaching WebAR application development; approaches to validating immersive technology teaching methodologies through quantitative and qualitative data triangulation. \textit{Further developed}: scientific propositions on the use of VR/AR technologies in education; approaches to ML integration in WebAR applications for educational purposes; approaches to qualitative analysis of pedagogical discourse in immersive technology education.

Practical significance: WebAR textbook (10 chapters), ``Synthetic Reality in Education'' course in LMS Moodle (18 modules), laboratory practicum (10 variants), virtual chemistry laboratory with augmented reality.

\textbf{Keywords}: immersive educational resources, WebAR, augmented reality, virtual reality, computer science teacher, professional training, pedagogical experiment, mixed methods, qualitative content analysis, CAMIL, TPACK, MindAR, Three.js, A-Frame, machine learning, TensorFlow.js.

\clearpage
% Список публікацій здобувача за темою дисертації (§4 Наказу МОН №40)
\printpratsi{}
